\chapter{Динамический захват и неединственность компакттонной орбиты}
\label{ch:v6-spectral-transition-discrete-compacton-dynamical-capture-gate}

\section{Почему существования и устойчивости недостаточно}

Предыдущие гейты нашли точную ветвь $\kappa=2\pi$, отсутствие расширяющей
флоке-моды и условное массово-размерное произведение. Для механизма
рождения требуется более сильное свойство: открытое множество обычных
локализованных состояний должно само приходить к одному физическому
endpoint или хотя бы оставлять в нём ненулевую устойчивую долю.

Литература по нелинейным квантовым блужданиям действительно содержит
солитонные режимы, но подчёркивает роль специально выбранных начальных
условий и нелинейной связи; см. \path{arXiv:1512.08358} и
\path{arXiv:2308.01014}. Поэтому точная орбита текущего проекта должна
пройти отдельный тест захвата, а не наследовать его из существования
солитонов в других моделях.

\section{Ретроспективная коррекция: точная орбита не изолирована}

Пусть
\begin{equation}
 v=(\ell,e)\in\mathbb C^2\oplus\mathbb C,
 \qquad
 \|\ell\|^2=|e|^2=\frac14.
 \label{eq:v6-capture-balanced-vector}
\end{equation}
Для составного хирального генератора выполняется
\begin{equation}
 K(v)v=\frac14v.
\end{equation}
При $\kappa=2\pi$ направление поэтому разворачивается на угол $\pi/2$.

Возьмём два произвольных, независимо выбранных вектора $v_L,v_R$,
удовлетворяющих \eqref{eq:v6-capture-balanced-vector}, и зададим
\begin{equation}
 \Psi_j=|{-}\rangle\otimes v_L,
 \qquad
 \Psi_{j+1}=|{+}\rangle\otimes v_R.
 \label{eq:v6-capture-general-two-site-family}
\end{equation}
После первого шага блоки меняются местами, а после второго возвращаются с
общим знаком:
\begin{equation}
 \boxed{F_{2\pi}^{,2}(\Psi)=-\Psi}.
 \label{eq:v6-capture-family-return-law}
\end{equation}

Таким образом, предыдущая ветвь $v_R=\pm i v_L$ действительно даёт
одношаговые собственные состояния с фазами $\pm i$, но не исчерпывает
локализованные решения. Любая относительная фаза даёт точный двухшаговый
компактон, а слабые направления и внутренние фазы $v_L,v_R$ могут
выбираться независимо.

Машинный сертификат проверил $32$ случайные точки этого семейства.
Максимальный остаток $\|F^2\Psi+\Psi\|$ равен
$2.14\cdot10^{-15}$. Скан шестнадцати относительных фаз дал ровно две
одношаговые собственные точки, $\varphi=\pi/2,3\pi/2$, но все шестнадцать
точек удовлетворяют \eqref{eq:v6-capture-family-return-law} с остатком не
выше $9.88\cdot10^{-16}$.

\begin{theorem}[Непрерывное компакттонное многообразие]
При $\kappa=2\pi$ минимальная составная монета имеет непрерывное семейство
точных двухузловых орбит $F^2=-1$. Ветви с собственной фазой $\pm i$
являются лишь специальным подсемейством $v_R=\pm i v_L$. Следовательно,
локализованный endpoint не единственен.
\end{theorem}

\section{Выбирает ли динамика баланс амплитуд}

Сначала изменена только доля нормы на левом узле:
\begin{equation}
 \|v_L\|^2=p,
 \qquad
 \|v_R\|^2=1-p,
\end{equation}
при сохранении хирального баланса внутри каждого блока. После $80$ шагов
получено:
\begin{center}
\begin{tabular}{@{}rrr@{}}
\toprule
$p$ & лучшая двухнаправленная вероятность & дефект многообразия \\
\midrule
$0.40$ & $0.502169$ & $0.705556$ \\
$0.49$ & $0.504576$ & $0.702036$ \\
$0.50$ & $1.000000$ & $1.29\cdot10^{-14}$ \\
$0.51$ & $0.504576$ & $0.702036$ \\
$0.60$ & $0.502169$ & $0.705556$ \\
\bottomrule
\end{tabular}
\end{center}
Точная точка $p=1/2$ сохраняется, но соседние состояния не возвращаются к
ней. Амплитудный баланс является резонансным условием, а не динамически
выбранной координатой.

\section{Выбирает ли динамика хиральный баланс}

Затем локальные нормы оставлены равными $1/2$, но доля слабой компоненты
заменена на $q$ вместо $1/2$. Результат после $80$ шагов:
\begin{center}
\begin{tabular}{@{}rrr@{}}
\toprule
$q$ & лучшая двухнаправленная вероятность & дефект многообразия \\
\midrule
$0.45$ & $0.918400$ & $0.105082$ \\
$0.48$ & $0.999850$ & $0.014141$ \\
$0.49$ & $0.999992$ & $0.007071$ \\
$0.50$ & $1.000000$ & $1.29\cdot10^{-14}$ \\
$0.51$ & $0.999992$ & $0.007071$ \\
$0.52$ & $0.999850$ & $0.014141$ \\
$0.55$ & $0.918400$ & $0.105082$ \\
\bottomrule
\end{tabular}
\end{center}
Малый хиральный дисбаланс может долго оставаться локализованным, но его
дефект не стремится к нулю. Динамика сохраняет окрестность семейства, не
выбирая точный сбалансированный слой.

\section{Проспективный тест обычных локализованных состояний}

До просмотра результата зафиксирован критерий захвата: на каждом
четырёхшаговом срезе от $t=60$ до $t=80$ дефект компакттонного
многообразия должен быть ниже $10^{-3}$, а вероятность в выбранной паре
направлений --- выше $0.99$. Для каждого из трёх ансамблей взято по
$12$ состояний с общим seed $20260821$:
\begin{center}
\begin{tabular}{@{}p{6.2cm}rrr@{}}
\toprule
Ансамбль & $N$ & захваты & $P_{\rm fin}^{\max}$ \\
\midrule
случайный один узел & 12 & 0 & $0.511118$ \\
случайные два узла & 12 & 0 & $0.518343$ \\
гауссов пакет со случайным внутренним вектором & 12 & 0 & $0.127962$ \\
\bottomrule
\end{tabular}
\end{center}
Наименьшие дефекты после времени прогрева равны соответственно
$0.675884$, $0.685747$ и $1.073676$, то есть ни один промах нельзя
объяснить только слишком строгим порогом $10^{-3}$.

\begin{proposition}[Захват в проверенном протоколе отсутствует]
Ни одно из $36$ проспективных обычных локализованных состояний не вошло в
компакттонное многообразие. Амплитудный и хиральный сканы также не показали
восстановления точных условий. Открытая область динамического захвата не
сертифицирована.
\end{proposition}

\begin{nogocriterion}
Численный скан не является доказательством отсутствия всех специальных
траекторий захвата. Но текущий компактон нельзя объявлять механизмом
рождения материи: он принадлежит непрерывному неединственному семейству,
а его баланс не выбирается динамикой и не достигается ни одним состоянием
заранее зафиксированных ансамблей.
\end{nogocriterion}

\section{Следующий гейт}

Следующий гейт
\path{version6_spectral_transition_discrete_compacton_branch_status_freeze_gate}
должен свести существование, устойчивость, масштаб и захват в единый
ledger и решить, остаётся ли неизменённая компакттонная архитектура
активной ветвью.

Нулевая гипотеза: без нового диссипативного, шумового или
неравновесного селектора ветвь сохраняется только как точное семейство
решёточных решений. Стоп-критерий: если R2, R3 и единственность R4 не
закрыты, дальнейшее добавление ручных селекторов запрещается, а ветвь
замораживается как математический результат.

\begin{protocol}
\begin{enumerate}
 \item точное семейство $F^2=-1$: \textbf{построено};
 \item независимые $v_L,v_R$: \textbf{допустимы};
 \item одношаговые фазы $\pm i$: \textbf{две специальные точки};
 \item изолированность endpoint: \textbf{нет};
 \item восстановление амплитудного баланса: \textbf{не найдено};
 \item восстановление хирального баланса: \textbf{не найдено};
 \item захваты в $36$ обычных испытаниях: \textbf{ноль};
 \item R2, эндогенный trigger: \textbf{не закрыт};
 \item R3, скорость: \textbf{не закрыт};
 \item R4, единственный устойчивый endpoint: \textbf{не закрыт}.
\end{enumerate}
\end{protocol}

Машинный сертификат сохранён в
\path{s2t_v6_spectral_transition_discrete_compacton_dynamical_capture_gate_results.json}.